\label{lecture05}
\subsection*{Цели и маршрут занятия (90 минут)}

Лекция связывает историю вычислительных систем с тем, что студент наблюдает
при запуске программы сегодня. После занятия студент сможет объяснить, зачем
операционной системе управлять совместным использованием компьютера; различить
исполняемый файл и процесс; описать на учебном уровне путь от команды в shell
до завершения процесса; объяснить назначение аргументов, окружения, трёх
стандартных потоков и кода завершения; диагностировать простые ошибки прав
доступа и запуска.

Примеры команд и демонстрация рассчитаны на Linux x86-64. Формат ELF и вызовы
Linux не являются общими для всех систем. Сравнение Windows и macOS ниже
показывает идеи, значимые программисту, а не исчерпывает их внутреннее
устройство.

\begin{center}
\small
\begin{tabular}{|p{0.73\linewidth}|r|}
\hline
Блок & Минуты \\
\hline
Вопрос: как сменить выполняемую работу? & 0--3 \\
Носитель и ручной запуск: неудобства, загрузчики, повторное использование & 3--15 \\
Пакетная очередь: автоматический переход между заданиями & 15--19 \\
От очереди к многопрограммности и разделению времени & 19--25 \\
Риск общей памяти: процесс, защита адресов и других ресурсов & 25--31 \\
Статус задания и сохранённый результат & 31--34 \\
Unix: стандартные потоки, перенаправление и pipe & 34--37 \\
Параллельные линии IBM mainframe и Unix; две модели запуска & 37--41 \\
Роли ОС: ресурсы, изоляция и режимы & 41--45 \\
От команды shell до нового процесса & 45--55 \\
Файл ELF, загрузчик и адресное пространство & 55--65 \\
Аргументы, окружение и демонстрация потоков & 65--79 \\
Системные вызовы, завершение и код возврата & 79--86 \\
Сравнение ОС и самопроверка & 86--90 \\
\hline
\end{tabular}
\end{center}

\subsection{Зачем нужна операционная система?}

Представим компьютер, на котором есть процессор, память, накопитель,
клавиатура и экран. Каждой программе пришлось бы самостоятельно договариваться
с остальными о том, кто и когда использует эти устройства. Если две программы
одновременно считают, где хранить их данные? Если одна программа ошиблась в
адресе, как не дать ей повредить работу другой? Если приложению нужно прочитать
файл, как связать имя файла с блоками на конкретном накопителе?

\textbf{Операционная система (ОС)} организует совместное использование
компьютера и предоставляет программам согласованный интерфейс. На примере
обычного компьютера её задачи включают:
\begin{itemize}
  \item создавать и учитывать выполняющиеся процессы, распределять процессорное
        время;
  \item предоставлять процессам память и ограничивать доступ одного процесса
        к памяти другого;
  \item организовывать файлы, каталоги и разрешения доступа;
  \item управлять устройствами через драйверы и общие интерфейсы;
  \item проверять полномочия пользователей и программ.
\end{itemize}

ОС не обязана быть видимым окном с кнопками. Графическая оболочка --- лишь один
из способов взаимодействовать с системой. В Linux команда в терминале проходит
через shell, но управление процессами, памятью и устройствами обеспечивает ОС.

\subsection{Историческая завязка: от ручного запуска к загрузчику}

Ранние электронные вычислительные машины не предлагали привычного рабочего
стола с независимыми приложениями. Оператор подготавливал программу и данные,
подавал их машине и начинал работу. Для ввода использовали перфокарты,
перфоленту, магнитную ленту и другие носители. Не следует представлять это как
универсальную схему «одна программа на одном носителе»: на носителе могли
храниться многие записи или задания. Трудность была в том, как выбрать нужную
работу, подать её машине и перейти к следующей без длительного ручного
переключения.

\subsubsection*{Какую работу снимал загрузчик?}

Машине нужно было прочитать представление программы, разместить его в памяти и
передать управление в нужную точку. Сначала этот механизм мог требовать
специальной процедуры или действий оператора. Затем загрузку поручали
небольшой программе-посреднику. Здесь важно различать начальную загрузку самой
системы (bootstrap) и загрузку очередной пользовательской программы: способы
первого старта машины различались, но после него ОС могла использовать общий
механизм для прикладных задач.

Следующее упрощение --- не писать отдельную последовательность загрузки для
каждой программы. Загрузчик мог находить разные программы, читать их и
передавать им управление. «Универсальный» в этом рассказе означает
\emph{повторно используемый для совместимых программ конкретной машины или
семейства}, а не способный исполнить любой код с любого носителя. Загрузчику
нужны согласованные сведения о формате, архитектуре и точке входа. Это всё ещё
не компиляция исходного текста: исходник C превращается в исполняемый файл
отдельно, что разбирается на лекции №7.

\subsection{Пакетные задания: очередь, результат, следующий шаг}

Когда оператор мог заранее подготовить несколько работ, стало выгодно запускать
их по очереди автоматически. Управляющая программа (monitor или supervisor)
читала очередь заданий и управляющие записи, загружала нужное приложение,
передавала ему данные и после окончания переходила дальше. Это и есть
\textbf{пакетная обработка} в рассматриваемой модели: серия заранее
подготовленных заданий исполняется без интерактивного диалога с пользователем.

Но автоматическая очередь сама по себе не делает процессор постоянно занятым.
Представим, что задание A после вычислений ждёт медленный ввод-вывод. Если
система допускает только последовательное выполнение, процессор простаивает
до завершения A, хотя следующее задание B уже готово.

Одно из решений --- \textbf{многопрограммная пакетная обработка}: несколько
заданий находятся в памяти, и пока A ожидает устройство, ОС выделяет процессор
готовому B. Устройство обслуживает A одновременно с вычислением B; на одном
процессоре в каждый момент исполняется только одно из них. Завершение
ввода-вывода переводит A в состояние готовности, но не обязательно немедленно
прерывает B. Очередь ожидающих запуска заданий отвечает на вопрос «что ещё
запустить?», а выбор среди готовых к исполнению --- «кому сейчас дать CPU?».
Пакетная система может поддерживать многопрограммность: это не противоположные
понятия.

Остаётся другая проблема: если B долго считает и не ждёт устройство,
переключения только при ожидании недостаточно для отзывчивой интерактивной
работы. В модели \textbf{разделения времени} сигнал таймера периодически
возвращает управление ядру. Планировщик решает, продолжать ли B или выделить
квант процессорного времени другой готовой задаче; состояние B сохраняется,
чтобы позже продолжить её работу. Таймер сам по себе не выбирает задачу, и не
каждое прерывание приводит к переключению. Эти решения объясняют переход от
очереди к многозадачности, но не описывают обязательную для всех ОС
историческую последовательность.

\subsubsection*{Вместе в памяти: почему нужна изоляция?}

Нескольким заданиям теперь нужно находиться в памяти и получать процессор
по очереди. Рассмотрим \emph{учебную машину без защиты памяти}: программа A
хранит данные по адресу \texttt{0x1200}, а программа B может записать по
этому же адресу. Ошибка указателя B --- или намеренная запись вредоносной
программы --- испортит данные A. Без разграничения B могла бы повредить и
код управляющей системы. Автоматическая очередь и переключение CPU сами по
себе не ограничивают доступ к памяти. Это иллюстрация проблемы, а не
утверждение, что все исторические пакетные машины были устроены одинаково.

Чтобы управлять независимыми работами, ОС учитывает \textbf{процесс}:
выполняющийся экземпляр программы с состоянием исполнения и связанными
ресурсами. Исполняемый файл остаётся на носителе; два запуска одного файла
создают два процесса с разным состоянием. Но само понятие процесса ещё не
гарантирует защиты: нужны аппаратно обеспеченные правила доступа.

В современной модели у каждого процесса есть \textbf{виртуальное адресное
пространство}. Например, адрес \texttt{0x1000} у A и B может вести к разным
физическим областям. ОС задаёт отображения и права, а процессор проверяет
их при обращении к памяти. Попытка обычной B записать в область A без
разрешения приводит к исключению, которое обрабатывает ОС. При этом ОС может
намеренно разрешить совместное использование страниц; изоляция не требует
запрета \emph{любого} обмена данными. Были и другие способы защиты памяти,
например проверка базы и границы: понятие процесса исторически не обязано
появляться вместе с виртуальной памятью.

Состояние процесса нужно для возобновления работы после переключения CPU;
к нему также привязаны открытые файлы и сведения о полномочиях. ОС отдельно
проверяет права доступа к файлам и устройствам. Виртуальный адрес не даёт
сам по себе права прочитать чужой файл или перенастроить устройство. Процессы
могут разделять ресурсы по правилам ОС: защита означает \emph{контролируемый}
доступ, а не обязательное владение каждым ресурсом в одиночку.

Чтобы не потерять результат и понять, как продолжать очередь, системе нужно
было сохранить сведения о завершении шага: штатно ли он закончился, возникла ли
ошибка, следует ли запускать следующий этап. В мэйнфреймовых ОС использовались
коды возврата и условия управляющего языка заданий (JCL), но точное значение
кода и область его действия зависели от системы. Код завершения Unix-процесса,
который shell позднее показывает через \texttt{\$?}, решает родственную задачу,
но это другой интерфейс, а не одно и то же историческое поле.

Результаты пакетной работы направляли в назначенные наборы данных, очереди
печати и журналы. Пользователь мог проверить отчёт после завершения запуска,
а отдельные диагностические сообщения помогали найти ошибку. Это создавало
важный рабочий цикл: \textbf{подготовить вход → выполнить задание → сохранить
результат и статус → решить, что запускать дальше}. Эта цепочка описывает
управление работами; при многопрограммности их исполнение может чередоваться.

\subsection{Unix: от сохранённого результата к композиции потоков}

В Bell Labs Unix начала формироваться в конце 1960-х годов; ранняя версия
работала в 1971 году, а значительную часть системы переписали на C в 1973.
Unix развивалась параллельно с IBM mainframe, а не после z/OS. Среди её идей
были интерактивная работа через shell, единообразный ввод-вывод файлов,
устройств и процессов и композиция небольших программ.

Записанный в файл результат можно передать следующей команде отдельным
запуском. Shell сделал соединение более непосредственным: стандартный ввод
\texttt{stdin}, обычный вывод \texttt{stdout} и поток диагностики
\texttt{stderr} можно было перенаправлять перед запуском. Операторы
\texttt{<}, \texttt{>} и \texttt{2>} связывают поток с файлом; канал
\texttt{|} связывает вывод одной программы со входом другой без промежуточного
файла. Таким образом, пакетный «результат → следующая работа» и Unix pipe
отвечают на близкую задачу передачи результата, но устроены по-разному:
очередь определяет порядок работ (их исполнение при многопрограммности может
перекрываться), а pipe непосредственно соединяет потоки процессов.

\subsubsection*{Кто управляет запуском: пакетный monitor и shell}

На двух схемах сопоставляются прежде всего \textbf{роли управления}. В
пакетном сценарии пользователь заранее передаёт системе очередь, а monitor как
управляющая часть этой системы выбирает следующий шаг и ведёт его до результата.
Здесь monitor --- исторический термин для управляющей программы; его не следует
автоматически отождествлять с ядром современной ОС.

В интерактивной Unix-подобной системе \textbf{shell сама является пользовательской
программой}. Она читает и разбирает команду, затем запрашивает у ядра запуск
внешней программы. Ядро создаёт процесс и загружает новый образ; shell обычно
ждёт завершения дочернего процесса и получает его статус. Встроенная команда
shell может выполняться без создания child-процесса. Поэтому различие схем не
сводится к выбору между \texttt{fork}/\texttt{exec} и каким-то другим
загрузчиком: одна показывает системное управление заранее заданной очередью,
другая --- пользовательский интерфейс, который для каждой внешней команды
обращается к ОС. Это различие сценариев, а не запрет пакетной обработки в Unix
или интерактивного доступа на мэйнфрейме.

\subsection*{Параллельные линии: мэйнфреймы IBM и Unix}

IBM объявила семейство System/360 и OS/360 в 1964 году; машины System/360
поступили в продажу в 1965 году. Мэйнфреймовая линия развивала управление
заданиями и ресурсами: позднее появились MVS (1970-е годы), а современным
представителем семейства является z/OS. В то же время в Bell Labs развивался
Unix: интерактивные процессы, файловый интерфейс, shell, перенаправления и
межпроцессные каналы.

Эти названия не составляют единую лестницу, пройденную всеми компьютерами.
Позднее массовые ПК дали ещё одну отдельную ветвь: PC DOS/MS-DOS, графическая
среда Windows поверх DOS, затем Windows NT с защищённой памятью и
многозадачностью. Современные Windows, Linux и macOS решают сходные задачи,
но их нативные форматы и программные интерфейсы различаются.

\begin{quote}
История привела нас к современному вопросу: когда пользователь просит
\texttt{./program}, кто проверяет файл, создаёт процесс и сообщает программе,
откуда читать данные и куда записывать результат?
\end{quote}

\subsection{Программа как файл и процесс как работа}

\textbf{Исполняемый файл} --- объект файловой системы с данными, которые ОС
может интерпретировать как программу для данной платформы. В Linux x86-64 таким
форматом часто является ELF (Executable and Linkable Format). В файле могут
находиться машинный код, данные, таблицы и сведения о том, как разместить
сегменты программы. Файл сам по себе не выполняет команды: он хранится на
накопителе.

\textbf{Процесс} --- выполняющийся экземпляр программы вместе с состоянием,
которым ОС управляет во время работы. У процесса есть виртуальное адресное
пространство, регистры и состояние исполнения, открытые ресурсы, аргументы
запуска и окружение. Один и тот же файл можно запустить несколько раз; тогда
существуют разные процессы с отдельным состоянием. Процесс может завершиться,
а исполняемый файл останется на диске.

Не путать эти понятия:
\begin{center}
\small
\begin{tabular}{|p{0.26\linewidth}|p{0.62\linewidth}|}
\hline
Что & Пример \\
\hline
Исходный код & \texttt{demo.c}, текст для компилятора \\
Исполняемый файл & \texttt{demo}, байты программы на диске \\
Процесс & Один работающий экземпляр \texttt{demo} \\
\hline
\end{tabular}
\end{center}

Подробное превращение исходного файла в исполняемый --- тема лекции №7. Здесь
начинаем с уже собранного файла и изучаем его запуск.

\subsection{Что происходит после команды запуска?}

Для команды:
\begin{lstlisting}
./demo first second
\end{lstlisting}
можно построить учебную последовательность:
\begin{enumerate}
  \item Shell разбирает введённую команду: имя программы, аргументы,
        перенаправления и прочие конструкции shell.
  \item Shell просит ОС запустить файл по указанному пути. В Unix-подобной
        системе shell обычно создаёт дочерний процесс, а затем этот процесс
        заменяет свой образ новой программой через семейство вызовов
        \texttt{exec}. Конкретные детали зависят от shell и способа запуска.
  \item ОС проверяет существование файла, его формат и право пользователя
        выполнить его. Она подготавливает новое виртуальное адресное
        пространство процесса.
  \item Загрузчик сопоставляет части исполняемого файла с областями памяти,
        настраивает начальное состояние и, при необходимости, подключает
        динамические библиотеки.
  \item Управление передаётся начальной точке программы. Для обычной программы
        на C до вызова функции \texttt{main} выполняется стартовый код среды C.
  \item При завершении процесса ОС сохраняет код завершения и освобождает
        принадлежащие ему ресурсы.
\end{enumerate}

Это не означает, что ОС читает весь файл и просто копирует его подряд в RAM.
Формат описывает области с разным назначением и правами; операционная система
и механизм виртуальной памяти могут загружать содержимое по мере необходимости.
Подробности динамической загрузки зависят от платформы и формата.

\subsubsection*{Изоляция и виртуальная память}

Виртуальное адресное пространство даёт каждому процессу собственную модель
адресов. Адрес \texttt{0x400000} в двух разных процессах не обязан обозначать
одну и ту же физическую ячейку памяти. Аппаратный механизм трансляции адресов
совместно с ОС отображает виртуальные области на физическую память или другие
ресурсы.

Это помогает изолировать процессы: обычная программа не может произвольно
прочитать память соседней программы. При недопустимом обращении процессор
сообщает о нарушении, а ОС обычно завершает процесс или передаёт управление
обработчику ошибки. Виртуальная память также позволяет ОС управлять общими
библиотеками, отображёнными файлами и страницами памяти.

\subsubsection*{Пользовательский режим, ядро и системные вызовы}

Обычная программа исполняется в \textbf{пользовательском режиме} и не может
самостоятельно выполнять привилегированные операции, например напрямую
перенастраивать устройство или таблицы памяти. Основная привилегированная часть
ОС называется \textbf{ядром} и работает в режиме ядра.

Чтобы попросить ОС выполнить защищённую операцию, программа использует
\textbf{системный вызов} (system call): например, запрос чтения, записи,
создания процесса или открытия файла. Переход в ядро --- контролируемая граница
с проверкой параметров и полномочий, а не обычный вызов функции C. Библиотечная
функция вроде \texttt{printf} может буферизовать данные и в итоге обратиться
к ОС для записи; не каждый вызов функции C означает отдельный системный вызов.

\subsection{Аргументы и окружение процесса}

В обычной программе C функция \texttt{main} часто записана как
\texttt{int main(int argc, char *argv[])}. Переменная \texttt{argc} содержит
число элементов, а \texttt{argv} --- массив строк. Элемент \texttt{argv[0]}
обычно содержит имя или путь, использованный для запуска; элементы с индексами
от 1 и выше --- аргументы пользователя. Аргументы передаются как текст: если
нужны числа, программа должна разобрать строки и проверить их.

\textbf{Окружение} --- набор пар «имя=значение», которые процесс получает при
запуске. Например, переменная \texttt{HOME} указывает домашний каталог, а
\texttt{PATH} содержит каталоги для поиска команд без явного пути. Процесс
может читать окружение через библиотеку, например функцией \texttt{getenv}.
Дочерняя программа обычно получает окружение от запускающей программы, но
родитель может передать изменённый набор. Секреты не следует считать
защищёнными только потому, что они находятся в переменной окружения.

\subsection{Три стандартных потока}

При запуске процесса в Unix-подобной среде обычно доступны три стандартных
потока:
\begin{itemize}
  \item \texttt{stdin} (дескриптор 0) --- источник ввода;
  \item \texttt{stdout} (дескриптор 1) --- обычный вывод программы;
  \item \texttt{stderr} (дескриптор 2) --- сообщения об ошибках и диагностика.
\end{itemize}

Поток не обязан быть физическим экраном или клавиатурой. Перед запуском shell
может направить его в файл или соединить с другой программой:
\begin{lstlisting}
./demo < input.txt > result.txt 2> errors.txt
./demo < input.txt | sort
\end{lstlisting}
В первой команде содержимое файла становится стандартным вводом; обычный вывод
и диагностика записываются в разные файлы. Во второй вывод первой программы
становится вводом \texttt{sort}. Символ \texttt{|} создаёт канал (pipe), по
которому команды обмениваются байтами.

\texttt{stdout} может быть буферизован: отдельные фрагменты вывода не обязательно
сразу видны на терминале или появляются в файле. Это объясняет, почему при
перенаправлении поведение вывода может отличаться по времени. Для разделения
данных и ошибок диагностические сообщения принято писать в \texttt{stderr}.

\subsection{Код завершения}

После окончания процесса ОС сообщает запускающей программе целочисленный
\textbf{код завершения}. По соглашению значение 0 обычно означает успех,
ненулевое значение --- некоторую ошибку или особый исход. Точное значение и
его смысл задаёт сама программа; универсальной таблицы кодов нет. В shell
\texttt{\$?} содержит код завершения последней выполненной foreground-команды:
\begin{lstlisting}
./demo
echo $?
\end{lstlisting}
В Unix-подобных системах наблюдаемое shell значение обычно ограничено диапазоном
от 0 до 255. В конвейере shell может показывать статус последней команды
конвейера, если не включён специальный режим вроде \texttt{pipefail}.

Код завершения --- не текстовое сообщение и не то же самое, что строка,
напечатанная в \texttt{stdout}. Программа может вывести сообщение об ошибке,
но случайно вернуть 0; поэтому вызывающей программе полезны оба канала:
текст для человека и статус для автоматического сценария.

\subsection{Права доступа и типовые ошибки запуска}

В Linux право чтения, записи и исполнения файла проверяется с учётом владельца,
группы и остальных пользователей, а также дополнительных механизмов защиты.
Команда \texttt{chmod +x demo} добавляет исполняемые биты для владельца,
группы и остальных пользователей, но не отменяет политики безопасности,
ограничения файловой системы или права родительских каталогов.

Имя команды без символа \texttt{/} обычно ищется в каталогах из переменной
\texttt{PATH}. Текущий каталог в \texttt{PATH} по умолчанию может отсутствовать,
поэтому для запуска файла из текущего каталога пишут \texttt{./demo}. Это
помогает случайно не запустить чужую программу с тем же именем.

\begin{center}
\small
\begin{tabular}{|p{0.32\linewidth}|p{0.58\linewidth}|}
\hline
Сообщение или симптом & Что проверить \\
\hline
\texttt{Permission denied} & Путь, право исполнения, права каталогов и ограничения системы \\
\texttt{command not found} & Имя команды и \texttt{PATH}; для файла в текущем каталоге --- \texttt{./demo} \\
\texttt{No such file or directory} & Опечатку в пути, файл и наличие нужного загрузчика/интерпретатора \\
\texttt{Exec format error} & Формат файла и соответствие архитектуре \\
\hline
\end{tabular}
\end{center}

Последний случай может удивить: файл существует, но ОС не может его исполнить.
Причиной бывает формат или архитектура файла; сообщение об отсутствующем файле
иногда относится к интерпретатору скрипта или динамическому загрузчику, а не
к самому исполняемому файлу. Команда \texttt{file demo} помогает определить
тип содержимого, но не гарантирует, что программу можно запустить.

\subsection{Linux, macOS и Windows: общий смысл и разные интерфейсы}

Во всех трёх системах программист встречает файлы программ, процессы, память,
ввод-вывод, права или политики безопасности и механизмы запуска. Но детали
различаются:
\begin{itemize}
  \item Linux обычно запускает нативные ELF-файлы; права исполнения и
        пользовательские/групповые права видны в интерфейсе файловой системы.
        Процессы и системные вызовы относятся к Linux, а не к Unix вообще.
  \item macOS основана на Darwin с компонентами BSD и XNU; нативные приложения
        часто распространяются как Mach-O и могут запускаться через графическую
        среду или shell. Модель разрешений и безопасности включает дополнительные
        механизмы платформы.
  \item Windows использует собственные форматы PE и интерфейсы создания
        процессов; программы запускают графически, через PowerShell, CMD или
        другие среды. WSL предоставляет Linux-среду, но не превращает обычный
        Linux ELF в нативную Windows-программу без соответствующего слоя.
\end{itemize}

Понятия «процесс», «аргументы», «стандартный ввод и вывод» похожи на уровне
задач программиста, но API, оболочки, форматы файлов и детали прав не следует
смешивать. Например, команда \texttt{chmod +x} относится к Unix-подобному
интерфейсу и не является универсальным способом включить запуск файла во всех
системах.

\subsection{Демонстрация и самопроверка}

В демонстрации используется программа \texttt{process-demo}. Она печатает
аргументы запуска, значение переменной \texttt{COURSE\_MESSAGE}, копирует
строки стандартного ввода в стандартный вывод и отдельно сообщает в
\texttt{stderr}. Это позволяет видеть, что потоками распоряжается не сама
программа в одиночку: shell может подключить к ним терминал, файл или канал.

Вопросы для обсуждения:
\begin{enumerate}
  \item Может ли один файл породить несколько процессов? Чем будут различаться
        эти процессы?
  \item Почему команда \texttt{demo} может не найти файл в текущем каталоге,
        а команда \texttt{./demo} найдёт?
  \item Куда попадёт текст, если перенаправить только \texttt{stdout}? А куда
        попадёт диагностика из \texttt{stderr}?
  \item Что сообщит \texttt{echo \$?}, если выполнить его после другой команды?
  \item Почему процесс не должен напрямую обращаться к устройству или памяти
        другого процесса?
  \item Почему история пакетных систем, Unix и Windows не является одной
        простой последовательностью поколений?
\end{enumerate}

Для наблюдения запущенных программ пригодны команды \texttt{ps} и \texttt{top}.
Дополнительный инструмент \texttt{strace} в Linux позволяет увидеть системные
вызовы процесса, но его вывод зависит от программы, версии системы и окружения.
Это средство исследования, а не перечень всех внутренних действий ОС.
